\begin{titlepage}

\projectAuthorName. \projectTitle. \projectType / supervisors \projectSupervisorNameFirst{} and \projectSupervisorNameSecond{}; \projectFaculty, Kaunas University of Technology.

Study field and area (study field group): \projectStudyFieldAndArea.

Keywords: \projectKeywords.

\projectCity, \projectYear. \pageref{LastPage}.

\begin{center}
  \textbf{Summary}
\end{center}

This literature analysis examines how functional near-infrared spectroscopy (fNIRS) and machine learning can support objective assessment of neurocognitive and mental-health related indicators. It does not treat the current evidence as sufficient for stand-alone diagnosis. A model trained on fNIRS data may classify task states, estimate workload, or detect patterns associated with a clinical group, but it becomes a diagnostic tool only when the data, labels, validation design, and clinical interpretation justify that claim.

The reviewed references are useful, but uneven. Several papers directly address fNIRS signal classification and deep learning methods, while others discuss fMRI-based autism classification or MRI segmentation. This analysis therefore frames the sources as methodological building blocks rather than as proof that fNIRS can already predict mental diseases with clinical reliability. The most defensible thesis position is exploratory: fNIRS can provide non-invasive physiological data; machine learning can learn patterns in those data; yet disease prediction requires stricter evidence than task classification.

\end{titlepage}






















